% Tabla 15: Contraste de hipótesis.
% Etiqueta original: tab:hipotesis.
% Fuente: capitulos/4-resultados.tex, lineas 235-248.
% Fragmento extraido automaticamente; requiere el preambulo de main.tex.

\begin{table}[htbp]
\centering\small
\caption{Contraste de hipótesis.}
\label{tab:hipotesis}
\begin{tabularx}{\textwidth}{@{}p{1.65cm} >{\raggedright\arraybackslash}X >{\raggedright\arraybackslash}X >{\raggedright\arraybackslash}p{2.4cm}@{}}
\toprule
Hipótesis & Criterio & Valor observado & Veredicto \\
\midrule
H1 & Prueba de signos exacta por réplica, $p < 0{,}05$, para C frente a A, \Aprime{} y B en LT & Frente a A y \Aprime: C mejor en 6 de 6 réplicas ($p = 0{,}031$) y en 24 de 24 combinaciones. Frente a B: C mejor en 3 de 6 réplicas ($p = 1$) y en 10 de 24 combinaciones & Se sostiene frente a A y \Aprime; no frente a B \\
H2 & Pendiente positiva de $\mathrm{MAE}_B - \mathrm{MAE}_C$ frente a la separación & Pendiente $-0{,}006$~mm/cm ($+0{,}09$; $-0{,}31$; $-0{,}28$; $-0{,}04$~mm) & No se sostiene \\
H3 & Cobertura $\geq 0{,}95$ y ancho predicho dentro de la mitad del observado & Cobertura 1,00 (72 de 72) en toda la pecera; ancho numérico 56--82~\%, razón predicho/numérico 1,006 & Se sostiene \\
\bottomrule
\end{tabularx}
\end{table}
